\begin{abstracteng}
\textit{Laser-Induced Breakdown Spectroscopy} (LIBS) offers rapid, multi-elemental analysis without sample preparation, making it highly suitable for geochemical characterization of volcanic soils. However, conventional quantification methods are limited by spectral complexity arising from overlapping emission lines, matrix effects, and plasma fluctuations. This research proposes a deep learning framework for quantitative elemental analysis of Mount Sinabung volcanic soil using a CNN--Transformer encoder--decoder architecture augmented with a Cross-Attention mechanism. Synthetic LIBS spectra are generated via a physics-based plasma model grounded in the Saha--Boltzmann equations under the assumption of Local Thermodynamic Equilibrium (LTE), utilizing atomic transition parameters from the NIST Atomic Spectra Database (ASD). The Cross-Attention module employs NIST line embeddings as queries to guide the spectral decomposition from polyatomic into per-element monoatomic components. To bridge the distribution gap between synthetic training data and real field measurements, domain adaptation is implemented via a Gradient Reversal Layer (GRL). Model performance is evaluated against X-ray Fluorescence (XRF) measurements as ground truth, targeting $R^2 > 0.95$ and MAPE $< 10\%$ across major rock-forming elements (Si, Al, Fe, Ca, Mg, K, Na, Ti). This approach is expected to contribute a scalable and interpretable solution for in-situ geochemical monitoring of volcanic activity in Indonesia.
}

\bigskip
\noindent
\textbf{Keywords:} LIBS, CNN-Transformer, Synthetic Spectra, Domain Adaptation, Cross-Attention, Volcanic Soil
\end{abstracteng}

